\section{Proofs and Boundary Cases}
\label{app:proofs}

\subsection{All-cutoff equivalence}

Fix a dated claim $c$ and a finite cutoff $\tau$.
If $\tau<t_c$, the claim has not appeared in either representation.
Otherwise define the witness set
\[
W_\tau(c)=\{d:t_c<t_d\leq\tau,\ G(d,c)\}.
\]
This set is nonempty exactly when its earliest possible timestamp satisfies $b_G(c)\leq\tau$.
Thus $W_\tau(c)=\varnothing$ exactly when $t_c\leq\tau<b_G(c)$.
This proves Theorem~\ref{thm:cutoff}, including the empty-witness case.
Equal timestamps never belong to $W_\tau(c)$.
The identifier tie rule changes the stored source, but not eligibility.

For a prefix $\mathcal C_{\leq\tau}$, the relevant witness set is exactly $W_\tau(c)$.
Claims outside the prefix cannot appear in it.
This proves Corollary~\ref{cor:prefix} without requiring correct or transitive pair decisions.
It requires fixed features and a fixed gate.
Corpus-dependent feature changes constitute different gates and belong to the influence analysis instead.

\subsection{Pair-local influence and temporal footprint}

Let $D=\{(d,c):t_d>t_c,\ G(d,c)\ne H(d,c)\}$.
Let $U=\{c:\exists d,\ (d,c)\in D\}$.
For each $c\notin U$, every possible witness has the same gate value under both compilers.
Therefore $b_G(c)=b_H(c)$.
Its certificate and eligibility agree at every cutoff.
This proves the mask containment and the bound $|U|\leq|D|=h$.

For the context bound, change the eligibility bits in $U$ one at a time.
Adding one eligible record either leaves the selected context unchanged or inserts it at its fixed ranking position.
At capacity, that insertion displaces one record.
Removing one eligible record either leaves the context unchanged or removes that record.
The next eligible record can fill its slot.
Each bit change therefore contributes at most two elements to the symmetric difference.
The triangle inequality gives $2|U|$.
Since each context has size at most $k$, the bound is also at most $2k$.

Both certificates start at $t_c$.
Their indicator functions differ exactly between their endpoints.
Hence Equation~\ref{eq:footprint} follows, including an infinite endpoint.
By linearity of expectation,
\[
\mathbb E_{\tau\sim\mu}|A_G(\tau)\triangle A_H(\tau)|
=\sum_c\mu(J_c).
\]
The context expectation is at most twice this quantity.
The statement uses the same ranking at each compared cutoff.
It does not require independent claim errors or independent query times.

\subsection{Amplification through a component bridge}

Use $m\geq2$ records with value $a$ at times $0,1,\ldots,m-1$.
Their shared attribute changes to $b$ at time $2m$.
A second attribute contains $z$ at time $m$ and $w$ at time $2m+1$.
Assume $a,b,z,w$ are distinct.
Initially all matching decisions are correct.
Each $a$ record has a replacement link to $b$.
The $z$ record has a replacement link to $w$.

Add the single false decision that $w$ replaces $b$.
The component graph now joins both histories.
Sorting this component makes $z$ the first later distinct value for every $a$ record.
Their component endpoints become $m$.
At cutoff $m+\frac12$, all $m$ records therefore become ineligible.
The direct construction changes only $b$'s boundary.
Because $b$ starts after the cutoff, direct eligibility does not change there.

The prefix at this cutoff contains neither endpoint of the false bridge.
Consequently, component reconstruction on that prefix cannot make the same exclusions.
This gives a prefix inconsistency caused entirely by later-dated records.
For a context difference, add $k$ independent filler records and rank the $a$ records first.
When $k\leq m$, direct selection returns $k$ such records.
Component selection replaces them with the fillers.
The contexts have symmetric difference $2k$.
This construction establishes possibility, not an observed frequency.

\subsection{Exact indexed compilation}

Order timestamp batches increasingly.
Within a batch, order candidate witnesses by identifier.
Before processing the first batch, the unresolved index is empty.
Assume it contains precisely earlier claims without an accepted earlier witness before some batch.
Complete posting lists expose every accepted pair from this batch to these claims.
The first positive comparison therefore has the earliest possible timestamp and the smallest eligible witness identifier.
Recording it establishes the exhaustive boundary and witness.
Removing the resolved claim is safe because subsequent batches have later timestamps.

After processing all comparisons, insert the batch's claims.
They cannot have a strictly later witness among previously processed records.
They therefore satisfy the unresolved invariant.
Induction establishes the output for every batch.
Any remaining claim has no accepted later witness and receives an infinite boundary.

For the supplied positive frame threshold, an accepted pair with nonempty frames shares at least one token.
The union of frame-token postings therefore contains every such candidate.
Empty frames use a common posting.
Repeated posting visits are deduplicated before the gate test.
For other gate definitions, completeness must be proved separately or ensured through a full unresolved scan.

\section{Source and Scoring Details}
\label{app:source}

The complete TempLAMA audit recovers all 300 keys used by the released adaptation.
It checks 2,677 original yearly records.
There are 613 multi-answer records and 245 keys with at least one such record.
The original builder takes the first answer identifier and collapses consecutive repetitions.
Rebuilding that procedure exactly reproduces the released claim texts and answer identifiers.
The audit finds 153 transitions where the previous selected answer remains in the next complete answer set.

Our source-set evaluation accepts all answer identifiers from the latest source observation at the cutoff.
It retains the existing claim pool and every method's previously computed ranking.
No parameter changes follow this rescore.
The old selected-answer oracle is not a complete source-set oracle.
It is therefore excluded from primary source-set superiority claims.

The single-answer subset contains keys with exactly one answer identifier in every source year.
This criterion selects 55 keys before examining retrieval outcomes.
The experiment retains all 747 claims as candidates.
Its 57 historical probes and 55 current probes therefore preserve the original distractor pool.

Historical cluster intervals resample keys and retain every sampled key's probes.
We use 5,000 resamples with seed 20261006.
This preserves dependence between dates from the same chronology.
All method comparisons use aligned questions and cutoffs.
The saved prediction records contain selected identifiers, target identifiers, and per-query metrics.

\section{Archived Current-Answer Runs}
\label{app:legacy}

The inherited reader uses \texttt{Qwen2.5-3B-Instruct}, greedy decoding, and at most 24 new tokens.
Its contexts contain five claims from a hybrid lexical and dense ranking.
The dense encoder is \texttt{all-MiniLM-L6-v2}~\citep{reimers2019sbert}.
Prompts include the available dates and request one answer value.
These settings and generated outputs precede the certificate study.
The compiler checks establish identical current masks under the unchanged pair gate.
Thus the archived contexts remain applicable to the current path.
They do not measure new historical reader performance.
\begin{table}[t]
\centering
\small
\begin{tabular}{lrrr}
\toprule
Corpus & $n$ & Flat & Original mask \\
\midrule
Wikidata & 33 & 42.42 & 96.97 \\
TempLAMA & 300 & 56.67 & 93.67 \\
Wikipedia prose & 35 & 71.43 & 85.71 \\
\bottomrule
\end{tabular}
\caption{Archived Qwen contained-value accuracy, in percent. These runs use the original single-target labels and ranking. They are inherited results, not fresh model evaluations.}
\label{tab:legacy-reader}
\end{table}


The archived scorer checks whether the normalized selected answer appears in the generated response.
We call this \emph{contained-value accuracy}.
A stricter version also rejects responses containing a competing archived answer.
Neither metric is normalized whole-answer exact match.
The TempLAMA rows use the original selected-answer transition task.
They do not use the complete source-answer scoring introduced in Section~\ref{sec:protocol}.

On that task, archived lifecycle selection reaches 0.937 contained-value accuracy across 300 TempLAMA questions.
Flat retrieval reaches 0.567 and the dated-prompt variant reaches 0.623.
On 33 Wikidata questions, the corresponding lifecycle and flat results are 0.970 and 0.424.
On 35 RealProse questions, they are 0.857 and 0.714.
The RealProse paired difference has archived interval $[0.000,0.286]$ and reported $p=0.076$.
These results describe the previous current-answer system.

The archived temporal-score comparisons adapt TempRALM and half-life formulas to the same local encoder.
They do not reproduce those papers' complete systems.
The selected half-life setting uses nine Wikidata development keys.
Those keys remain in the 33-key archived aggregate.
The full result cards preserve that selection history.

\section{Implementation and Reproduction}
\label{app:reproduction}

The released compiler uses the Python standard library and the supplied feature parser.
Its output records claim identifiers, timestamps, witness identifiers, source hashes, and detector settings.
It also records pair tests and posting visits.
An infinite endpoint is encoded as a null endpoint in JSON.
An undated start has no historical eligibility and remains separately identifiable.

The correctness runner compares indexed output with an exhaustive gate scan.
It checks repeated values, shuffled inputs, equal timestamps, missing timestamps, and arbitrary pair relations.
Constructed cases are labeled as software checks.
Benchmark results come from the supplied source collections.
The run manifest distinguishes new CPU experiments from inherited model generations.
The experiment scripts preserve per-query and per-intervention outputs alongside aggregate tables.

Each historical certificate identifies the source that caused exclusion.
Each evaluation record identifies the claims actually selected.
These records support inspection of an interval error separately from a ranking or reader error.

\section{AI Assistance}
AI tools assisted with literature searches, code development, proof drafts, and manuscript editing.
No new language-model inference contributed to the revised experimental tables.
